% 雨果·斯坦豪斯（Hugo Steinhaus）（综述）
% license CCBYSA3
% type Wiki

本文根据 CC-BY-SA 协议转载翻译自维基百科\href{https://en.wikipedia.org/wiki/Hugo_Steinhaus}{相关文章}。

\begin{figure}[ht]
\centering
\includegraphics[width=6cm]{./figures/2db84a4d97edc48c.png}
\caption{雨果·斯坦豪斯（1968年）} \label{fig_YGstgs_1}
\end{figure}
雨果·迪奥尼济·斯坦豪斯（Hugo Dyonizy Steinhaus，英语发音：/ˈhjuːɡoʊ ˈstaɪnˌhaʊs/，波兰语发音：[ˈxuɡɔ ˈʂtajnxaws]，1887年1月14日－1972年2月25日）是一位波兰数学家和教育家。斯坦豪斯于1911年在哥廷根大学师从大卫·希尔伯特（David Hilbert）获得博士学位，后来成为利沃夫扬·卡齐米日大学（现今乌克兰利维夫）的教授，并在那里帮助建立了后来被称为“利沃夫数学学派”（Lwów School of Mathematics）的学术中心。他被誉为“发现”数学家斯特凡·巴拿赫（Stefan Banach）的人，与巴拿赫共同在泛函分析领域作出了重要贡献，特别是提出了著名的巴拿赫–斯坦豪斯定理（Banach–Steinhaus theorem）。

第二次世界大战后，斯坦豪斯在弗罗茨瓦夫大学（Wrocław University）数学系的建立以及战后波兰数学界的重建中发挥了重要作用。斯坦豪斯一生发表了约170篇科学论文与著作，他在数学的多个分支领域留下了深远影响，包括泛函分析、几何学、数理逻辑以及三角学。值得注意的是，他被认为是博弈论（game theory）和概率论（probability theory）的早期奠基者之一，为后续学者发展更系统的理论框架奠定了基础。
\subsection{早年生活与求学经历}
斯坦豪斯于1887年1月14日出生于奥匈帝国的亚斯沃（Jasło）一个有犹太血统的家庭。\(^\text{[1]}\)他的父亲博古斯瓦夫（Bogusław）是一名地方实业家，拥有一家砖厂，同时也是一位商人；他的母亲名为埃韦丽娜（Ewelina），娘家姓利普希茨（Lipschitz）。斯坦豪斯的叔叔伊格纳奇·斯坦豪斯（Ignacy Steinhaus）\(^\text{[2]}\)是一名律师，也是“波兰议会俱乐部”（Koło Polskie）的成员，并担任加利西亚和洛多梅里亚王国（Kingdom of Galicia and Lodomeria）地区议会——加利西亚议会（Galician Diet）的代表。
\begin{figure}[ht]
\centering
\includegraphics[width=8cm]{./figures/844c35d7f6c20eff.png}
\caption{1907年，斯坦豪斯于哥廷根（Göttingen）拍摄的合影中站在最右侧。} \label{fig_YGstgs_2}
\end{figure}
1905年，斯坦豪斯在亚斯沃文理中学（gymnasium in Jasło）完成学业。他的家人原本希望他成为一名工程师，但他却被抽象数学深深吸引，开始自学当时著名数学家的著作。同年，他进入伦贝格大学（今乌克兰利维夫大学，University of Lemberg）学习哲学与数学。\(^\text{[2]}\) 1906年，他转入哥廷根大学（Göttingen University）。\(^\text{[1]}\)在那里，他师从大卫·希尔伯特（David Hilbert），并于1911年获得博士学位，论文题目为《Dirichlet 原理的新应用》（德文原题：Neue Anwendungen des Dirichlet'schen Prinzips）。\(^\text{[3]}\)

第一次世界大战爆发后，斯坦豪斯回到波兰，加入约瑟夫·毕苏斯基（Józef Piłsudski）领导的波兰军团（Polish Legion）服役，战后定居于克拉科夫（Kraków）。\(^\text{[4]}\)斯坦豪斯是无神论者。\(^\text{[5]}\)
\subsection{学术生涯}
\subsubsection{两次世界大战之间的波兰时期}
在1916年至1917年期间，即波兰于1918年重新获得独立之前，斯坦豪斯曾在克拉科夫（Kraków）为内务部工作，服务于短暂存在的傀儡政权——波兰王国（Kingdom of Poland）。\(^\text{[6]}\)

1917年，他开始在伦贝格大学（University of Lemberg，后来的波兰扬·卡齐米日大学 Jan Kazimierz University）任教，并于1920年取得讲师资格（habilitation）。\(^\text{[6]}\) 1921年，他被任命为“副教授”（profesor nadzwyczajny），1925年晋升为“正教授”（profesor zwyczajny）。\(^\text{[6]}\)在此期间，他开设了当时仍属前沿领域的勒贝格积分理论（Lebesgue integration）课程——这是法国以外最早开设此类课程的大学之一。\(^\text{[3]}\)

在利沃夫（Lwów）任教期间，斯坦豪斯共同创立了著名的“利沃夫数学学派”（Lwów School of Mathematics）。\(^\text{[7]}\)他也是苏格兰咖啡馆（Scottish Café）中活跃的数学家圈成员之一，不过据斯坦尼斯瓦夫·乌拉姆（Stanisław Ulam）回忆，斯坦豪斯本人在这些聚会上通常更喜欢到街上那家更为高雅的茶馆小坐。\(^\text{[4]}\)
\subsubsection{第二次世界大战时期}
\begin{figure}[ht]
\centering
\includegraphics[width=6cm]{./figures/6c566e9ff83fe38a.png}
\caption{《苏格兰笔记本》（The Scottish Book）——利沃夫数学学派（Lwów School of Mathematics）的著名记录集，斯坦豪斯曾为其作出贡献，并极有可能在第二次世界大战期间将其保存了下来。} \label{fig_YGstgs_3}
\end{figure}
1939年9月，纳粹德国与苏联分别入侵并占领波兰，以履行先前签署的《莫洛托夫–里宾特洛甫条约》。利沃夫（Lwów）最初落入苏联控制之下。斯坦豪斯曾考虑逃往匈牙利，但最终决定留在利沃夫。苏联方面对大学进行了重组，使其更具乌克兰特色，但他们任命了斯坦豪斯的学生斯特凡·巴拿赫（Stefan Banach）为数学系主任，斯坦豪斯也得以恢复教学工作。数学系师资队伍此时还因多位来自德占波兰的难民学者而得到加强。斯坦豪斯后来写道，在此期间的经历使他对“一切苏联官员、政客和委员们都产生了不可克服的生理性厌恶”。\(^\text{[A]}\)

在两次大战之间以及苏联占领期间，斯坦豪斯向著名的《苏格兰笔记本》（Scottish Book）贡献了十个问题，其中包括最后一个——记录于1941年纳粹发动“巴巴罗萨行动”（Operation Barbarossa）攻占利沃夫前夕。\(^\text{[4]}\)

由于犹太血统，斯坦豪斯在纳粹占领期间被迫躲藏。他最初藏身于利沃夫的朋友家中，随后辗转至扎莫希奇（Zamość）附近的小镇奥西奇纳（Osiczyna），以及克拉科夫（Kraków）附近的贝尔德霍夫（Berdechów）。\(^\text{[7][8]}\) 波兰反纳粹抵抗组织为他提供了一份伪造的身份证明，身份是早已去世的森林护林员“格热戈日·克罗赫马尔内”（Grzegorz Krochmalny）。以此身份，他秘密教授课程（德军占领期间禁止波兰人接受高等教育）。  由于担心一旦被捕便会被处死，斯坦豪斯在无书可查的情况下，仅凭记忆重建并记录下了自己掌握的全部数学知识，并撰写了大量回忆录——其中只有少部分在战后被发表。\(^\text{[8]}\)

此外，在藏匿期间，由于与外界隔绝、无法获得可靠的战争信息，斯坦豪斯设计了一种统计方法，通过分析地方报纸上零星刊登的讣告，来估算德军在前线的伤亡人数。  
该方法基于讣告中“死者为某人之子”、“某人之二子”、“某人之三子”等表述的相对频率。\(^\text{[8]}\)

据他的学生兼传记作者马克·卡茨（Mark Kac）回忆，斯坦豪斯曾说，自己一生中最幸福的一天，是“德军撤离波兰而苏军尚未来临的那二十四小时”——他说：“他们已经走了，而他们还没来。”\(^\text{[8]}\)
\subsubsection{第二次世界大战之后}
\begin{figure}[ht]
\centering
\includegraphics[width=6cm]{./figures/6e60811ae1628b9a.png}
\caption{纪念牌，波兰弗罗茨瓦夫（Wrocław）} \label{fig_YGstgs_4}
\end{figure}
在第二次世界大战的最后几天，斯坦豪斯仍然处于隐匿状态中。他听到一个传言：利沃夫大学（University of Lwów）将被迁往布雷斯劳（Breslau，即今日波兰的弗罗茨瓦夫 Wrocław），这一城市将在《波茨坦协定》（Potsdam Agreement）后划归波兰（而利沃夫则被并入苏维埃乌克兰）。虽然起初他对此消息持怀疑态度，但他拒绝了来自乌兹（Łódź）和卢布林（Lublin）的教职邀请，最终前往弗罗茨瓦夫，并在那里开始于弗罗茨瓦夫大学（University of Wrocław）任教。\(^\text{[7]}\)在弗罗茨瓦夫期间，斯坦豪斯重启了他在利沃夫“苏格兰笔记本”（Scottish Book）中的旧理念——即让杰出或有潜力的数学家们记录他们感兴趣的问题，并附上奖励方案。他创办了《新苏格兰笔记本》（New Scottish Book），以延续那一传统。  极有可能正是斯坦豪斯本人在战争期间保存了原版《苏格兰笔记本》，并在战后将其寄给斯坦尼斯瓦夫·乌拉姆（Stanisław Ulam），由后者译成英文。\(^\text{[4]}\)

在斯坦豪斯的努力下，弗罗茨瓦夫大学逐渐发展成为波兰数学重镇，正如战前的利沃夫大学一样享誉学界。\(^\text{[8]}\)

1960年代，斯坦豪斯曾担任美国圣母大学（University of Notre Dame，1961–1962年）及英国萨塞克斯大学（University of Sussex，1966年）的访问教授。\(^\text{[4][9]}\)
\subsection{数学贡献}
斯坦豪斯一生共发表了约170篇学术论文与著作。\(^\text{[4]}\)不同于其学生斯特凡·巴拿赫（Stefan Banach）在泛函分析领域的专攻，斯坦豪斯的研究范围极为广泛，涵盖几何学、概率论、泛函分析、三角与傅里叶级数理论，以及数理逻辑。\(^\text{[3][4]}\)  他还涉足应用数学，并热衷与工程师、地质学家、经济学家、医生、生物学家，甚至用卡茨（Mark Kac）的话说，“连律师都合作过”。\(^\text{[8]}\)

他在泛函分析中的最重要贡献之一，是与巴拿赫共同于1927年提出并证明的Banach–Steinhaus 定理，这一结果现已成为该领域的基础工具之一。

斯坦豪斯对博弈问题的兴趣，使他早期就提出了“策略”的形式化定义，这一思想比约翰·冯·诺依曼（John von Neumann）系统化的博弈论表述还要早数年。因此，他被视为现代博弈论的早期奠基人之一。\(^\text{[6]}\)在研究无限博弈问题的过程中，他与另一位学生扬·米切尔斯基（Jan Mycielski）共同提出了决定性公理（Axiom of Determinacy）。\(^\text{[8]}\)

斯坦豪斯同时也是概率论的早期创立者与奠基人之一。当时的概率论尚处萌芽阶段，甚至尚未被视为数学的正式分支。\(^\text{[8]}\)他首次以测度论的公理体系形式刻画了“掷硬币”问题，这一思想对十年后俄国数学家安德烈·柯尔莫哥洛夫（Andrey Kolmogorov）建立概率论公理体系产生了直接影响。\(^\text{[8]}\)他还首次严格定义了“事件相互独立”的含义，以及“随机变量服从均匀分布”的定义。\(^\text{[4]}\)

在二战藏匿期间，斯坦豪斯研究了著名的公平分蛋糕问题（fair cake-cutting problem）：  即如何将一个异质资源在不同偏好的多人之间分配，使每个人都认为自己得到了“比例公平”的份额。  斯坦豪斯的工作被认为开启了现代公平分配问题的研究方向。\(^\text{[B]}\)

此外，斯坦豪斯也是第一个提出火腿三明治定理（ham-sandwich theorem）猜想的人之一，\(^\text{[8][10]}\)并且是最早提出k-均值聚类方法（k-means clustering）思想的数学家之一。\(^\text{[11]}\)
\subsection{遗产与影响}
据说，斯坦豪斯在1916年“发现”了波兰数学家斯特凡·巴拿赫（Stefan Banach）。当时，他在克拉科夫（Kraków）的一座公园里偶然听到有人提到“勒贝格积分”（Lebesgue integral）一词，于是上前攀谈。后来他称巴拿赫是自己“最伟大的数学发现”。\(^\text{[12]}\) 此后，斯坦豪斯与巴拿赫及另一位当时在场的讨论者奥托·尼科迪姆（Otto Nikodym）共同创立了克拉科夫数学会（Mathematical Society of Kraków），该学会后来演变为波兰数学会（Polish Mathematical Society）。\(^\text{[4]}\)他是多家学术机构的成员，包括波兰科学院（PAN）、波兰学术院（PAU）、波兰数学会（PTM）、弗罗茨瓦夫科学学会（Wrocław Scientific Society, WTN）[pl]，以及多家国际科学团体和科学院。\(^\text{[6]}\)

1921年，斯坦豪斯在《数学基础》（Fundamenta Mathematicae）上发表了首批论文之一。\(^\text{[13]}\)1929年，他与巴拿赫共同创办了《数学研究》（Studia Mathematica），\(^\text{[7]}\)1953年又创立了《数学的应用》（Zastosowania Matematyki），并参与创建了《数学会刊》（Colloquium Mathematicum）与《数学专论》（Monografie Matematyczne）等重要刊物。\(^\text{[1]}\)

他先后获授多所大学的名誉博士学位，包括：华沙大学（Warsaw University, 1958）；弗罗茨瓦夫医学院（Wrocław Medical Academy, 1961）；波兹南大学（Poznań University, 1963）；弗罗茨瓦夫大学（Wrocław University, 1965）。\(^\text{[14]}\)

斯坦豪斯精通多种外语，以其格言和哲思著称。他的一本收录波兰语、法语与拉丁语格言的小册子在他逝世后出版。\(^\text{[8]}\)

2002年，波兰科学院与弗罗茨瓦夫大学联合举办了纪念活动——“2002：雨果·斯坦豪斯年”（2002, The Year of Hugo Steinhaus），以表彰他对波兰及世界科学的贡献。\(^\text{[15]}\)

数学家马克·卡茨（Mark Kac），亦是斯坦豪斯的学生，在缅怀恩师时写道：

“他是那所奇迹般在两次世界大战之间于波兰盛开的数学学派的奠基人之一。  
或许正是他，比任何人都更大程度地，使波兰数学从二战的灰烬中重生，并恢复到如今受人尊敬的强盛地位。  
他是一位有着深厚文化修养的人，从最好的意义上来说，他是西方文明的产物。”\(^\text{[3]}\)
\subsection{主要著作}
\begin{itemize}
  \item 《什么是数学，什么不是数学》}（Czym jest, a czym nie jest matematyka, 1923）\(^\text{[14]}\)
  \item 《奇点凝聚原理》（Sur le principe de la condensation de la singularités，与巴拿赫合著，1927）\(^\text{[3]}\)
  \item 《正交级数理论》（Theorie der Orthogonalreihen，与斯特凡·卡茨马日 Stefan Kaczmarz 合著，1935）\(^\text{[3][16]}\)
  \item 《数学万花筒》（Kalejdoskop matematyczny，英文版题为 Mathematical Snapshots，1939）\(^\text{[3][14]}\)
  \item 《弗罗茨瓦夫分类学》（Taksonomia wrocławska，与他人合著，1951）
  \item 《有限点集的连接与分割》（Sur la liaison et la division des points d'un ensemble fini，与他人合著，1951）\(^\text{[17]}\)
  （该文是对博鲁夫卡算法（Borůvka's algorithm）的多次独立再发现之一）
  \item 《一百道初等数学题》（Sto zadań: One Hundred Problems in Elementary Mathematics，1964）\(^\text{[4][18]}\)
  \item 《正面还是反面》（Orzeł czy reszka: Heads or Tails，1961）\(^\text{[19]}\)
  \item 《理性词典》（Słownik racjonalny: A Rational Dictionary，1980）\(^\text{[20]}\)
\end{itemize}
\subsection{家庭}
他的女儿莉迪娅·斯坦豪斯（Lidya Steinhaus）嫁给了波兰文学评论家兼戏剧理论家扬·科特（Jan Kott）。
\subsection{参见}
\begin{itemize}
  \item 弗雷林对称公理（Freiling's Axiom of Symmetry）
  \item 七分之一面积三角形（One-seventh area triangle）
  \item 约翰逊–特罗特算法（Johnson–Trotter Algorithm）
  \item 斯坦豪斯猜想（Steinhaus Conjecture）
  \item 斯坦豪斯多边形记号（Steinhaus Polygon Notation）
  \item 斯坦豪斯定理（Steinhaus Theorem）
  \item 斯坦豪斯测长器（Steinhaus Longimeter）
  \item 最后减量者法（Last Diminisher, 公平分割算法）
\end{itemize}
\subsection{注释}
\begin{enumerate}[label=\Alph*.]
  \item “我对一切苏联官员、政治家与委员产生了无法克服的生理性厌恶。”（摘自杜达著作，第23页）
  \item 两人版本的“公平分蛋糕问题”可由经典儿童规则“一个人分，另一个人选”解决。  
  斯坦豪斯是第一个将该问题推广至三人或更多参与者的人，并提出了“比例公平”（proportional division）标准。
\end{enumerate}
\subsection{参考文献}
\begin{enumerate}
  \item 《一百道初等数学题》序言，Courier Dover Publications，1974年，第4页，ISBN 978-0-486-23875-3。
  \item 雅斯沃镇（Jasło）官方网站（2010）。〈Steinhaus Hugo Dyonizy〉，收录于“Mieszkaniec: Steinhaus Hugo Dyonizy. Jasło. Moje miasto, nasz wspólny dom.”，2011年10月1日存档，检索于2011年8月16日。
  \item Kac, Mark（1974）。〈Hugo Steinhaus——回忆与致敬〉（PDF）。《美国数学月刊》（The American Mathematical Monthly），第81卷第6期，美国数学学会（Mathematical Association of America），页572–581。doi:10.2307/2319205。JSTOR 2319205。2011年9月27日从原PDF存档。
  \item John O'Connor 与 Edmund F. Robertson（2000年2月）。〈Hugo Dyonizy Steinhaus〉，《MacTutor 数学史档案》。苏格兰圣安德鲁斯大学数学与统计学院，检索于2011年8月16日。
  \item Steven G. Krantz（2002）。《数学轶事录：数学家的故事与趣闻》（Mathematical Apocrypha: Stories and Anecdotes of Mathematicians and the Mathematical）。美国数学学会出版，第202页。ISBN 978-0-88385-539-3。作者指出斯坦豪斯是一位直言不讳的无神论者。
  \item Monika Śliwa（2010年5月4日）。〈Hugo Steinhaus〉，弗罗茨瓦夫大学官网。2011年10月5日存档。
  \item Duda, Roman（2005）。〈战后弗罗茨瓦夫数学的起点〉（Początki Matematyki w Powojennym Wrocławiu）（PDF）。《大学评论》（Przegląd Uniwersytecki），9月刊。波兰数学会弗罗茨瓦夫分会。2011年9月27日存档，检索于2011年8月17日。
  \item Kac, Mark（1987）。《偶然的谜：自传》（Enigmas of Chance: An Autobiography）。加州大学出版社，第49–53页。ISBN 978-0-520-05986-3。
  \item Chełminiak, Wiesław（2002）。〈欧洲的弗罗茨瓦夫〉（Wrocław Europy）。刊于《Wprost》。检索于2011年8月20日。
  \item Beyer, W. A. 与 Andrew Zardecki（2004）。〈火腿三明治定理的早期历史〉（The Early History of the Ham Sandwich Theorem）。《美国数学月刊》，第111卷第1期，页58–61。doi:10.2307/4145019。JSTOR 4145019。ProQuest 203746537。
  \item Lindsten, Frederik; Ohlsson, Frederik; Lennart, Ljung（2011）。〈放松并开始聚类：$k$ 均值聚类的凸化〉（Just Relax and Come Clustering: A Convexification of k-means Clustering）。瑞典林雪平大学自动控制技术报告，第1页。
  \item Feferman, Anita Burdman 与 Feferman, Solomon（2004）。《阿尔弗雷德·塔斯基：生命与逻辑》（Alfred Tarski: Life and Logic）。剑桥大学出版社，第29页。ISBN 978-0-521-80240-6。
  \item Kuratowski, Kazimierz 与 Borsuk, Karol（1978）。〈《数学基础》第一百卷〉（One Hundred Volumes of Fundamenta Mathematicae）（PDF）。《Fundamenta Mathematicae》，第100卷，波兰科学院，第3页。
  \item 《雨果·斯坦豪斯教授》（Prof. Hugo Steinhaus），弗罗茨瓦夫理工大学官网。
  \item Aleksander Weron（2002年1月4日）。〈2002——雨果·斯坦豪斯年〉（2002 – Rok Hugona Steinhausa）。弗罗茨瓦夫理工大学，检索于2011年8月26日。
  \item Stefan Kaczmarz 与 Hugo Steinhaus（1951）。《正交级数理论》（Theorie der Orthogonalreihen）。切尔西出版社（Chelsea Pub. Co.），检索于2011年9月2日。
  \item Steinhaus 等（1951）。〈关于有限点集的连接与分割〉（Sur la liaison et la division des points d'un ensemble fini）（PDF）。波兰虚拟科学图书馆——数学文集。
  \item Steinhaus, Hugo（1974）。《一百道初等数学题》（One Hundred Problems in Elementary Mathematics）。Courier Dover Publications。ISBN 978-0-486-23875-3。
  \item Steinhaus, Hugo（1961）。《正面还是反面》（Orzeł czy reszka）（波兰语），第1卷，华沙：Państwowe Wydawnictwo，OCLC 68678009。
  \item Steinhaus, Hugo（1980）。《理性词典》（Słownik racjonalny）（波兰语），第1卷，奥索林斯基国家出版社（Zakład Narodowy im. Ossolińskich），OCLC 7272718。
\end{enumerate}
\subsubsection{延伸阅读}
\begin{itemize}
  \item 卡齐米日·库拉托夫斯基（Kazimierz Kuratowski），《波兰数学五十年：回忆与反思》（A Half Century of Polish Mathematics: Remembrances and Reflections），牛津：Pergamon Press，1980年，ISBN 978-0-08-023046-7，第173–179页及其他相关章节。
  \item 雨果·斯坦豪斯（Hugo Steinhaus），《数学万花筒》（Mathematical Snapshots），第二版，牛津，1951年，书背简介（blurb）。
\end{itemize}
\subsection{外部链接}
\begin{itemize}
  \item \href{https://www.mathgenealogy.org/id.php?id=8016}{数学谱系计划（Mathematics Genealogy Project）上的雨果·斯坦豪斯（Hugo Steinhaus）条目}
  \item \href{https://mathscinet.ams.org/mathscinet/search/author.html?mrauthid=70027}{MathSciNet 上的雨果·斯坦豪斯（Hugo Steinhaus）}
  \item \href{https://zbmath.org/authors/?q=ai:steinhaus.hugo}{Zentralblatt MATH（数学文摘）上的雨果·斯坦豪斯（Hugo Steinhaus）}
\end{itemize}
